\documentclass[UTF8,fontset=none,11pt,a4paper]{ctexart}
\usepackage{ipara-handbook}
\handbooksetup{DS-A04 贪心选择与交换证明}{408 · 数据结构算法 Extension}{Local Choice · Exchange · Frontier}
\begin{document}
\handbooktitle{贪心选择与交换证明}{局部最优何时可以安全地成为全局承诺}{408 · 数据结构 Algorithm Extension A04}

\begin{corebox}{母模型：选择、承诺、交换与剩余问题}
贪心不是“每一步选最大的”，而是
\[
\text{Feasible Family}
\to\text{Candidate Ordering}
\to\text{Safe Choice Proof}
\to\text{Residual Problem}
\to\text{Repeat}.
\]
安全选择意味着：存在一个全局最优解包含该选择，或贪心前缀在某个序关系下始终不落后于任意竞争方案。没有交换、cut、stays-ahead、matroid 或其他等价证明时，局部规则只是启发式。
\end{corebox}

\begin{methodbox}{Owner 与停止边界}
\begin{itemize}
\item \kw{Owns：}局部承诺的证明义务、区间调度/合并、跳跃可达边界。
\item \kw{Uses：}DS11 排序、DS05 优先队列、DS06 Union-Find、DS08 图上的 cut/relax；DS-A02 的上界剪枝。
\item \kw{Stops：}参与证明的具体排序、堆、并查集和图算法机制回各自 Owner；无法证明的策略不稳定纳管。
\end{itemize}
\end{methodbox}
\tableofcontents\newpage

\section{区间调度：最早结束留下最多尾部空间}

要选择最多个互不重叠区间，按右端点递增排序；扫描时选择第一个与已选区间兼容的区间。设任意最优解的第一个区间为 $I$，贪心选择 $G$ 的右端点不晚于 $I$。把 $I$ 替换为 $G$ 后，后续任意与 $I$ 不重叠的区间也与 $G$ 不重叠，因此替换不减少可选数量。重复该交换即可得到包含贪心首选的最优解。

LC 435 的“最少移除”先转成“最多保留互不重叠”；答案是总数减最大保留数。LC 452 的气球是区间刺点：按右端点选择一支箭，所有与当前点相交的区间被同一箭覆盖；遇到左端点大于当前点才需要新箭。闭区间与开区间的相等端点规则必须固定，`>=` 与 `>` 不能凭模板互换。

端点合同必须先写在草稿上：闭区间中 $[a,b]$ 与 $[b,c]$ 通常相交，兼容条件是新左端点 $>end$；半开区间 $[a,b)$ 中相接可共存，兼容条件是新左端点 $\ge end$。气球刺点把区间视为闭区间交集，判定下一个左端点 $\le$ 当前箭点时仍可共用一箭。空输入、单区间、相同右端点和完全包含区间不能靠默认初始化掩盖；若端点已排序，扫描是 $O(n)$，一般输入的排序才是 $O(n\log n)$ 主成本。

\subsection{合并区间不是同一个贪心目标}

区间合并按左端点排序，维护当前合并段的最右边界；若下一区间左端点不超过当前右端点，就扩张右端点，否则提交一个已闭合段。它不在选择“最多区间”，而是在排序后维护覆盖包络；正确性来自已处理前缀的合并结果与未处理区间的分离，不应套用右端点调度的证明。

区间列表交集是另一种双指针接口：两列表各自按左端点有序，当前交集为 $[\max(l_1,l_2),\min(r_1,r_2)]$，若交集非空则输出，然后推进右端点较小的一侧。它利用的是“较早结束者不可能再与未来区间形成更长当前交集”，不是把两列表先拼成一个集合。视频拼接、最少箭头和区间合并也不能因为都出现区间就共享同一排序键。

\begin{warnbox}{区间题的三种不同问题}
最大互不重叠、最少点覆盖全部区间、合并所有重叠段看似相似，排序键和不变量却不同。先写目标的集合语义，再决定右端点选择、左端点扫描或其他方法；“区间题按右端点排序”不是普遍规则。
\end{warnbox}

\section{跳跃游戏：stays-ahead 而非真的跳到最远点}

LC 55 维护扫描前缀能到达的最远位置 `farthest`。循环不变量是：扫描到 $i$ 时，所有 $\le i$ 的位置若未超过 `farthest` 都可达；若出现 $i>farthest$，后续位置更不可能到达。更新 $farthest=\max(farthest,i+nums[i])$，最终检查是否覆盖末端。这是可达前缀的支配关系，不需要构造具体跳法。

空数组和单元素数组应先按题目合同处理；单元素通常已在终点。若中间出现零，只有当当前 `farthest` 无法越过该位置时才构成失败；扫描到终点后可以立即停止。若要求恢复一条路径，仅有 `farthest` 不足，还需保存使边界改进的前驱或在等价最短层中二次回溯。

LC 45 求最少跳数时，`[0,curEnd]` 是当前跳数可到达的一层，扫描该层时维护下一层最远边界 `nextEnd`；到达 `curEnd` 才把跳数加一并提交 `nextEnd`。算法并不是每次从当前位置真的跳到最远下标，而是证明任何下一跳都不可能越过该层的最大覆盖边界，因此一次层扩张等价于 BFS 的最少边数。

扫描层边界时，只有在到达 `curEnd` 才支付一次跳数；在层内选择某个位置并不会立即增加代价。若输入不保证可达，应在提交前检查 `nextEnd\le curEnd`，否则会在不可达零区间上伪造进展。该算法把同一层的可达集合压成一个区间，若转移产生的集合不是连续可覆盖区间，就不能直接套用此边界压缩。

\section{证明工具箱}

\subsection{Exchange：把最优解换成贪心解}

证明格式：取任意最优解 $O$；若它已经包含贪心选择 $G$，继续处理剩余问题；否则找到 $O$ 中与 $G$ 竞争的位置，用 $G$ 替换该位置，并证明可行性与目标值不变或不差。区间调度的“更早结束”就是把尾部空间变大。

\subsection{Stays-ahead：每个前缀都不落后}

对跳跃游戏、扫描覆盖等问题，未必存在一个固定的“包含贪心选择”的最优解，但可证明贪心在每个扫描前缀的覆盖范围至少与任意方案一样大。于是任何方案能到达的后续位置，贪心也不会更晚到达。

Stays-ahead 证明不等于“每次选数值最大的跳”。它比较的是同一资源预算下的\kw{整个前缀覆盖边界}；只要一个候选在当前位置大却导致后续边界不占优，就不能被局部数值替代。证明失败的启发式只能暂存为 Candidate，不进入 Canonical 贪心规则。

\subsection{Cut / Relax：图算法的接口}

Kruskal 的最小跨切边可安全加入，依赖 cut property；Prim 每次从已选集合跨出的最小边也用同一性质。Dijkstra 选择当前最小 tentative distance 后 settle，依赖所有边权非负；其“局部最小”不是一般贪心，而是带有 DS08 的 relax 不变量和边权前提。图算法的机制正文不在本册复制，本册只抽取证明接口。

\section{堆辅助的贪心选择}

当候选不是一次性排序，而是动态产生时，DS05 优先队列可维护当前最小/最大候选。例如合并多个有序流、调度任务或 Prim/Dijkstra，每次从堆顶取“当前安全候选”，再把新邻接候选放回。堆只提供取极值的局部能力，安全性仍由调用算法的证明提供；不能因为用了堆就自动得到贪心正确性。

\section{贪心与 DP、回溯的分流}

若局部选择一旦做出就能通过交换证明安全，贪心可以丢弃其他分支；若不同历史必须比较未来最优值，则需要 DP；若必须保留所有路径、剪枝只能依赖约束或上界，则使用回溯/分支限界。一个问题可以同时有贪心和 DP 解法，但要分别说明哪种结构性质使贪心更快，以及它在哪些输入条件下失效。

\begin{examplebox}{跳跃游戏的 DP 对照}
可以定义 $F(i)$ 表示到达位置 $i$ 的最少跳数，得到 $O(n^2)$ 的 DP；贪心把所有同一跳数可达的位置合并成一个 frontier，并只保留下一层最远边界，利用 stays-ahead 证明把状态压成 $O(1)$。这不是“贪心总比 DP 好”，而是该问题的可达集合具有区间闭包。
\end{examplebox}

\section{复杂度、失败边界与测试}

排序驱动的区间算法通常为 $O(n\log n)$ 时间；已排序输入或线性扫描部分才是 $O(n)$。跳跃游戏为 $O(n)$、$O(1)$ 额外空间，但若要恢复具体跳跃路径，需要保存 parent 或在等价边界内二次选择。动态候选用堆时成本增加为每次 $O(\log k)$。

攻击测试包括：区间端点相等、空区间、包含关系、全重叠、全不重叠、跳跃数组含零、不可达末端、多个并列最优、负权图和负数区间。若把非负权换成负权、把可交换区间换成带权区间或加入资源约束，原贪心可能立即失效，应转交 DS08/DS-A03 或建立新的证明。

对区间调度可用小规模全枚举对拍最大非重叠数量；对合并区间检查输出覆盖原并集且结果互不重叠；对 Jump II 用 BFS 层数作为基线，验证边界恰好到达、零步长断点和不可达输入。若需要输出方案而非只输出数量，额外检查 parent/箭点选择是否与最优值证明相容。

\section{压缩复原}

忘记贪心规则时问：候选按什么量排序？选择后剩余问题是否同型？能否把任意最优解交换成包含该选择的解？或能否证明每个扫描前缀不落后？这个证明失败时，停止背口诀，不要把“看起来合理”写成算法正确性。

\begin{corebox}{一句话压缩}
贪心的核心不是局部最优，而是局部选择带有可验证的全局承诺；排序、堆和指针只负责产生候选，exchange、stays-ahead 或 cut property 才负责证明它不会错。
\end{corebox}
\end{document}
